\documentclass[12pt,a4paper]{article}
\usepackage[utf8]{inputenc}
\usepackage[indonesian]{babel}
\usepackage{amsmath}
\usepackage{amsfonts}
\usepackage{amssymb}
\usepackage{geometry}

\geometry{margin=2.5cm}

\title{Menghitung Volume dan Luas Permukaan dengan Teorema Pappus}
\author{Ahmad Fakhrul Bawani}
\date{}

\begin{document}

\maketitle

\section*{Soal}

The square region with vertices $(11,0)$, $(22,11)$, $(11,22)$, and $(0,11)$ is revolved about the $x$-axis to generate a solid. Find the volume and surface area of the solid.

\subsection*{1. Karakteristik Daerah Persegi}
Mari kita tentukan terlebih dahulu panjang sisi, luas, keliling, dan pusat massa (sentroid) dari daerah persegi tersebut:
\begin{itemize}
  \item \textbf{Panjang Sisi ($s$):}
    Menggunakan rumus jarak antara titik $(11,0)$ dan $(0,11)$:
    \begin{equation*}
      s = \sqrt{(11 - 0)^2 + (0 - 11)^2} = \sqrt{121 + 121} = \sqrt{242} = 11\sqrt{2}
    \end{equation*}
  \item \textbf{Luas Persegi ($A$):}
    \begin{equation*}
      A = s^2 = (11\sqrt{2})^2 = 242
    \end{equation*}
  \item \textbf{Keliling Persegi ($L$):}
    \begin{equation*}
      L = 4s = 4(11\sqrt{2}) = 44\sqrt{2}
    \end{equation*}
  \item \textbf{Sentroid Persegi ($\bar{x}, \bar{y}$):}
    Pusat massa persegi terletak pada titik potong diagonalnya, yaitu rata-rata dari koordinat keempat titik sudutnya:
    \begin{equation*}
      \bar{y} = \frac{0 + 11 + 22 + 11}{4} = \frac{44}{4} = 11
    \end{equation*}
    *(Kita hanya memerlukan koordinat $\bar{y}$ karena poros putaran adalah sumbu-$x$, sehingga jarak sentroid ke poros adalah $R = \bar{y} = 11$).*
\end{itemize}

\subsection*{2. Menghitung Volume Benda Putar ($V$)}
Berdasarkan **Teorema Pappus untuk Volume**, volume benda putar yang dihasilkan dari perputaran suatu daerah sejauh $360^\circ$ mengelilingi poros yang tidak memotongnya adalah:
\begin{equation*}
  V = 2\pi \cdot R \cdot A
\end{equation*}

Substitusikan nilai jarak sentroid $R = 11$ dan luas daerah $A = 242$:
\begin{align*}
  V &= 2\pi \cdot 11 \cdot 242 \\
  V &= 22\pi \cdot 242 \\
  V &= 5324\pi
\end{align*}

Jadi, volume benda putar tersebut adalah \textbf{$5324\pi$ satuan volume}.

\subsection*{3. Menghitung Luas Permukaan Benda Putar ($S$)}
Berdasarkan **Teorema Pappus untuk Luas Permukaan**, luas kulit dari benda putar yang dihasilkan oleh busur/kurva pembatas luar berpanjang $L$ adalah:
\begin{equation*}
  S = 2\pi \cdot R \cdot L
\end{equation*}

Substitusikan nilai jarak sentroid kurva keliling $R = 11$ dan keliling persegi $L = 44\sqrt{2}$:
\begin{align*}
  S &= 2\pi \cdot 11 \cdot 44\sqrt{2} \\
  S &= 22\pi \cdot 44\sqrt{2} \\
  S &= 968\pi\sqrt{2}
\end{align*}

Jadi, luas permukaan benda putar tersebut adalah \textbf{$968\pi\sqrt{2}$ satuan luas}.

\end{document}